\documentclass[11pt]{ctexart}
\usepackage[a4paper,margin=1in]{geometry}
\usepackage{graphicx}
\usepackage{xcolor}
\usepackage{hyperref}
\definecolor{refgray}{gray}{0.45}
\title{\textbf{市场最新观点汇总}\\[0.4em]{\large 全球资管增长逻辑切换、AI从试点走向规模部署、新兴市场改革与通胀压力并存}}
\author{KC桌面 自动生成}
\date{260703}
\begin{document}
\maketitle
\tableofcontents
\newpage
\section{一页摘要}
\begin{itemize}
  \item 全球资管行业超80\%收入增长依赖市场上涨，净流入能力成为唯一护城河，行业进入'增长不增收'阶段。
  \item 银行AI部署卡在'试点'与'规模'之间，区域银行反而领先于大型银行，组织灵活性成为关键瓶颈。
  \item IMF对新兴市场（吉尔吉斯斯坦、柬埔寨）的改革评估显示，增长引擎切换与通胀压力并存，结构性改革窗口正在收窄。
  \item 代币化不会消灭金融基础设施，但将重新定义功能分配，混合治理模型成为最可能终局。
  \item 物理AI投资回报周期已从5-7年缩短至1-3年，产业决策者面临'必须行动'的临界点。
\end{itemize}
\section{投行/券商}
今日暂无新增报告。

\begin{itemize}
  \item 今日暂无新增报告。
\end{itemize}
\section{战略咨询}
波士顿咨询与麦肯锡共同揭示：资管行业增长逻辑从规模转向净流入，AI部署从试点走向规模的关键在于组织架构而非技术，物理AI正跨越投资临界点。

\subsection{资管行业：增长幻觉终结，净流入能力成为唯一护城河}
\textbf{全球资管规模突破147万亿美元，但超80\%收入增长来自市场上涨，行业利润率15年未变（约30\%），成本增速（5.4\%）已超过收入增速（5.1\%），'负经营杠杆'成为常态。}

\begin{itemize}
  \item 2024-2025年全球资管行业收入增长的80\%以上由市场估值上升驱动，若市场转向震荡或下行，行业收入基础将严重侵蚀。
  \item 机构管理费以每年3\%的速度下降，被动产品和ETF持续主导净流入，每新增一美元管理规模带来的平均管理费持续降低。
  \item 零售客户贡献了2020-2025年全球资管规模增长的61\%，亚太地区增速最快（年增9\%），但传统资管在数字平台和社交媒体渠道几乎无存在感。
  \item 美国被动基金前10家公司拿走了超90\%的净流入，私募市场前50家PE机构2024年募资占比37\%（远超十年均值22\%），渠道和平台成为新的'看门人'。
  \item BCG建议资管机构从'产品能力'转向'渠道嵌入能力'和'技术基础设施能力'，AI部署需先理顺分销四大支柱，否则无法产生效果。
\end{itemize}
\begin{figure}[htbp]
\centering
\includegraphics[width=0.88\linewidth]{figures/fig_100.jpg}
\caption{EXHIBIT 1 - This will require asset managers to operate differently. Firms that align with the next sources of capital and adapt how they compete will capture a disproportionate share of future growth. \#\# EXHIBIT 1 \# Market Performance Drove Over 80\% of Gross Revenue Grow}
\end{figure}
\begin{figure}[htbp]
\centering
\includegraphics[width=0.88\linewidth]{figures/fig_102.jpg}
\caption{Exhibit 3 - Global AuM has more than tripled and revenue more than doubled over the past 15 years. Yet industry profit margins remain close to 30\%, roughly where they stood in 2010. Between 2010 and 2025, revenues grew at 5.1\% annually while costs rose slightly faster at }
\end{figure}
\begin{figure}[htbp]
\centering
\includegraphics[width=0.88\linewidth]{figures/fig_103.jpg}
\caption{Exhibit 4 - The rise of digital-native investors is concentrating flows in a smaller set of platforms that act as gatekeepers to capital. For asset managers, success will depend on being embedded in these ecosystems, with products and capabilities designed for how capital}
\end{figure}
{\small\color{refgray}\textbf{References:} [R007] 波士顿咨询：资管行业增长幻觉终结，净流入能力成为唯一护城河}

\subsection{AI部署：银行'谨慎'正在变成竞争劣势，组织架构比算法更关键}
\textbf{麦肯锡对44家全球金融机构调研显示，52\%已将生成式AI列为首要任务，但实际落地进展缓慢，区域银行在用例部署数量上反而领先于超大型银行，组织灵活性成为胜负手。}

\begin{itemize}
  \item 银行AI应用能力分为三类：摘要能力、内容生成能力、客户交互能力，目前大多数机构进展集中在低风险的'摘要'领域，如早期预警系统和信贷决策支持。
  \item 区域银行（资产1000-5000亿美元）在AI用例部署数量上领先于超大型银行，反映出组织灵活性比规模更重要。
  \item BCG将个性化决策科学分为三层：倾向性与uplift评分、上下文老虎机、基于基础模型的智能体推理，但绝大多数企业连第一层都没做好，20\%-40\%营销预算被浪费。
  \item 营销组织核心单元必须从'战役'切换为'货架'，成熟模式下70\%-80\%的客户触点将从预定义营销战役转向AI代理实时调取的个性化交互。
  \item 银行面临的真正挑战不是技术成熟度，而是组织在战略层面将AI从'效率工具'升级为'竞争杠杆'的意愿和能力，ROI难以早期量化是主要障碍。
\end{itemize}
\begin{figure}[htbp]
\centering
\includegraphics[width=0.88\linewidth]{figures/fig_122.jpg}
\caption{Exhibit 2 - Commitment to implementation of gen AI, by type of institution, \$\textasciicircum{}\{1\}\$ number Adoption is a priority Interested, but not a clear priority Not a priority McKinsey \& Company Gen AI offers financial institutions three highly useful capabilities: concision, meanin}
\end{figure}
\begin{figure}[htbp]
\centering
\includegraphics[width=0.88\linewidth]{figures/fig_126.jpg}
\caption{Exhibit 4 - Somewhat surprisingly, the group most advanced in deployment is regional banks, which are ahead of megabanks in number of use cases (Exhibit 4). In addition, core regionals are most advanced on ideation and planning. Very few use cases have reached the stage o}
\end{figure}
\begin{figure}[htbp]
\centering
\includegraphics[width=0.88\linewidth]{figures/fig_119.jpg}
\caption{Exhibit 1 - NBA action selection occurs across three layers: propensity and uplift scoring, contextual bandits, and agent-based reasoning. These layers are additive, not alternative. Systems don't necessarily need to graduate from one to the next, however, as each layer h}
\end{figure}
{\small\color{refgray}\textbf{References:} [R008] 波士顿咨询：营销的终极形态，不是“旅程”而是“原子动作”; [R010] 波士顿咨询：营销战役的终结，不是营销的终结; [R011] 波士顿咨询：下一代个性化决策，关键不在算法而在组织架构; [R012] 麦肯锡：银行对AI的“谨慎”正在变成一种竞争劣势}

\subsection{物理AI：回报周期从5-7年缩至1-3年，产业决策者进入'必须行动'阶段}
\textbf{由于计算机视觉、虚拟训练环境和软件定义架构的突破，物理AI的自动化工作范围扩大50\%，工程时间减少70\%，投资回报周期从5-7年缩短到1-3年，已跨越从'值得关注'到'必须行动'的临界点。}

\begin{itemize}
  \item 传统机器人约75\%的成本来自系统集成、安装和工程调试，而软件定义的物理AI可将这部分成本削减超过一半。
  \item 1-3年的回报周期意味着，即使一家企业现在才开始规划物理AI部署，在大多数CEO的常规任期之内就能看到财务影响。
  \item 自动化不再只是资本密集型巨头的游戏，能够快速决策、灵活试错的中型企业也获得了参与竞争的资格。
  \item BCG建议CEO重新评估运营效率、整体设计工作流、先想好技术架构再找供应商，并建立内部AI能力中心。
\end{itemize}
\begin{figure}[htbp]
\centering
\includegraphics[width=0.88\linewidth]{figures/fig_117.png}
\caption{波士顿咨询视觉摘要 1 - 波士顿咨询｜波士顿咨询：BCG：Physical AI的回报周期已从5-7年缩至1-3年，CEO们还在等什么？｜用于快速识别该外部报告的核心问题和市场含义。}
\end{figure}
{\small\color{refgray}\textbf{References:} [R009] 波士顿咨询：BCG：Physical AI的回报周期已从5-7年缩至1-3年，CEO们还在等什么？}

\section{智库/国际机构}
IMF多份报告揭示：新兴市场增长引擎切换与通胀压力并存，代币化正重新定义金融基础设施功能分配，SDR新分配暂缓但留有活口。

\subsection{新兴市场：增长红利消退，改革窗口收窄}
\textbf{IMF对吉尔吉斯斯坦和柬埔寨的评估显示，过去几年的高增长驱动力（转口贸易、侨汇、基建支出）正在接近天花板，通胀攀升、财政转赤字、信贷过快增长成为共同挑战，结构性改革紧迫性上升。}

\begin{itemize}
  \item 吉尔吉斯斯坦2021-2025年实际GDP累计扩张47\%，但2026年3月通胀已攀升至11\%（远超央行5-7\%目标区间），财政从连续三年盈余转为赤字，信贷增速接近50\%。
  \item 吉尔吉斯斯坦增长引擎正在从贸易红利切换至基建驱动，但中间存在空档期——2025年出口已下降14.5\%，转口贸易进入平台期。
  \item 柬埔寨预算执行改革取得显著进展，但FMIS系统'最后一公里'、现金垫款灰色地带、国库单一账户实质性落地仍是关键短板，IMF建议2027年前将所有未授权实体转化为完全授权实体。
  \item 吉尔吉斯斯坦真实货币政策远比政策利率宽松：银行体系结构性流动性过剩导致利率走廊'地板'取代政策利率成为市场利率实际锚，信贷扩张正从外部成本推动转向国内需求拉动通胀。
  \item IMF建议吉尔吉斯斯坦维持数据依赖的紧缩立场、增加汇率弹性、推进国企改革和减少非正规经济，中小企业占GDP约50\%但融资受限。
\end{itemize}
\begin{figure}[htbp]
\centering
\includegraphics[width=0.88\linewidth]{figures/fig_006.jpg}
\caption{Figure 1 - 42. Staff recommends that the next Article IV consultation be held on a standard 12-month cycle. Sources: country authorities and IMF staff estimates. \#\# Figure 1. Kyrgyz Republic: Real Sector and Social Indicators Growth outpaced comparators for the past 4 ye}
\end{figure}
\begin{figure}[htbp]
\centering
\includegraphics[width=0.88\linewidth]{figures/fig_007.jpg}
\caption{Figure 1 - \#\# Figure 1. Kyrgyz Republic: Real Sector and Social Indicators Growth outpaced comparators for the past 4 years,... Inflation is above target... ...thanks to trade and construction. ...threatening recent poverty reduction. Poverty Headcount Ratio at \textbackslash{}\$8.30 pe}
\end{figure}
\begin{figure}[htbp]
\centering
\includegraphics[width=0.88\linewidth]{figures/fig_056.png}
\caption{IMF视觉摘要 1 - IMF｜IMF：吉尔吉斯斯坦三年狂飙后，真正的考验是通胀与改革窗口｜用于快速识别该外部报告的核心问题和市场含义。}
\end{figure}
{\small\color{refgray}\textbf{References:} [R001] IMF报告：柬埔寨预算执行改革走到了哪一步，又卡在了哪里; [R003] IMF：吉尔吉斯斯坦三年狂飙后，真正的考验是通胀与改革窗口; [R004] IMF报告：吉尔吉斯共和国的真实货币政策远比政策利率宽松}

\subsection{代币化：混合治理模型成为最可能终局，兼容者将胜出}
\textbf{IMF两份工作论文指出，代币化不会消灭金融基础设施，但会重新分配功能——确定性、规则驱动的流程可迁移至智能合约，但需要法律确定性和主观判断的功能仍必须由机构承担，最可能的终局是'代码+机构'的混合模式。}

\begin{itemize}
  \item 可以迁移到链上的功能包括：资产发行和记录、所有权转移、结算中的DvP执行、抵押品估值和自动追缴、标准化数据报告。
  \item 不能完全迁移的功能集中在需要法律效力或主观判断的领域：新交易的法律确认、违约管理中的自由裁量、跨账本协调、监管特殊数据访问、压力情景下的保证金参数调整。
  \item 银行与加密企业正在'双向奔赴'：JPM Coin部署在Coinbase的无许可链Base上，UBS加入Tempo公共测试网，SG在以太坊、Solana、Stellar和XRP Ledger上发行欧元稳定币，但都通过白名单机制叠加'许可控制'。
  \item IMF提出三种架构模型（单一账本、兼容账本、共同账本），不同模型下风险分布完全不同——单一账本效率最高但单点风险集中，兼容账本灵活性好但跨链协调成本高。
  \item 基础设施层正在发生的'混合治理'实验将比资产层本身的产品创新更具决定性，谁能在开放性与合规性之间找到可持续平衡点，谁就能在下一阶段金融架构中占据枢纽位置。
\end{itemize}
\begin{figure}[htbp]
\centering
\includegraphics[width=0.88\linewidth]{figures/fig_002.jpg}
\caption{Figure 1 - Tokenization does not imply disintermediation, but institutional redesign. The most plausible outcome is a hybrid FMI model, in which smart contracts perform a greater share of operational and transactional functions, while legal entities remain responsible fo}
\end{figure}
\begin{figure}[htbp]
\centering
\includegraphics[width=0.88\linewidth]{figures/fig_003.jpg}
\caption{Figure 2 - \#\# Possible architectures The prior two sections discussed properties of blockchains in the abstract. In reality, the relationship between chains, assets, and owners – the architecture of the tokenized financial system – matters. Architecture encompasses the n}
\end{figure}
\begin{figure}[htbp]
\centering
\includegraphics[width=0.88\linewidth]{figures/fig_004.jpg}
\caption{Figure 3 - The first model, called the single ledger model, entails all owners having access to the same ledger on which the assets being transacted are recorded (Figure 3). Two parties might exchange a bond for money on the same ledger, for instance. The simplicity of t}
\end{figure}
{\small\color{refgray}\textbf{References:} [R002] IMF：Tokenization不会消灭金融基础设施，但会重新定义谁做什么; [R006] IMF报告：Tokenization正在重塑金融基础设施，但真正的赢家是“兼容者”}

\subsection{全球储备资产：SDR新分配暂缓，但危机触发条件仍存}
\textbf{IMF总裁明确在第十三个基本期内不会主动提议新一轮SDR普遍分配，因为当前全球经济不缺流动性和储备资产，且主要股东政治支持不到位，但报告留有活口——出现'意料之外的重大发展'时可重新提出分配提案。}

\begin{itemize}
  \item 自2021年SDR分配以来，全球储备资产增长了超过21\%，新兴和发展中国家（不含中国）增长了超过22\%，SDR存量相对于全球储备、资本流动和国际贸易的比例都显著高于2021年分配前的水平。
  \item 新SDR分配需要获得85\%投票权支持，IMF总裁与各执行董事沟通后，多数投票权代表明确表示不支持在当前背景下启动新分配。
  \item 历史经验显示，SDR分配需要经济危机级别的触发条件：2009年全球金融危机、2021年新冠危机是最近两次，再往前是1981年之后28年没有任何普遍分配。
  \item IMF对'长期全球储备资产补充需求'的定义决定了未来什么级别的危机才能重新打开SDR分配大门，报告中的图1和图2展示了全球储备资产、资本流动和贸易相对于SDR存量的历史变化。
\end{itemize}
\begin{figure}[htbp]
\centering
\includegraphics[width=0.88\linewidth]{figures/fig_089.jpg}
\caption{Figure 2 - Sovereign Issuances, 2005–2025 (In billions of USD) Figure 2. Long-term Perspective of SDR Holdings, 1980–2025 (In percent of) Global Reserves 1/ Global Gross Capital Flows Global Trade 3/ 1/ Including gold.}
\end{figure}
\begin{figure}[htbp]
\centering
\includegraphics[width=0.88\linewidth]{figures/fig_090.jpg}
\caption{Figure 2 - Figure 2. Long-term Perspective of SDR Holdings, 1980–2025 (In percent of) Global Reserves 1/ Global Gross Capital Flows Global Trade 3/ 1/ Including gold. 2/ Excluding China. 3/ Goods and services.}
\end{figure}
\begin{figure}[htbp]
\centering
\includegraphics[width=0.88\linewidth]{figures/fig_092.png}
\caption{IMF视觉摘要 1 - IMF｜IMF：SDR新分配的门已经关上，但钥匙还在桌上｜用于快速识别该外部报告的核心问题和市场含义。}
\end{figure}
{\small\color{refgray}\textbf{References:} [R005] IMF：SDR新分配的门已经关上，但钥匙还在桌上}

\section{结语}
今日国际信源主线清晰：资管行业增长逻辑从规模转向净流入，AI部署从试点走向规模的关键在于组织架构而非技术，新兴市场面临增长引擎切换与通胀压力并存的考验，代币化正重新定义金融基础设施功能分配。物理AI已跨越投资临界点，产业决策者需加速行动。
\newpage
\section{报告覆盖清单}
\begin{itemize}
  \item [R001] 智库/国际机构 | 智库/国际机构 / 新兴市场：增长红利消退，改革窗口收窄 - IMF报告：柬埔寨预算执行改革走到了哪一步，又卡在了哪里
  \item [R002] 智库/国际机构 | 智库/国际机构 / 代币化：混合治理模型成为最可能终局，兼容者将胜出 - IMF：Tokenization不会消灭金融基础设施，但会重新定义谁做什么
  \item [R003] 智库/国际机构 | 智库/国际机构 / 新兴市场：增长红利消退，改革窗口收窄 - IMF：吉尔吉斯斯坦三年狂飙后，真正的考验是通胀与改革窗口
  \item [R004] 智库/国际机构 | 智库/国际机构 / 新兴市场：增长红利消退，改革窗口收窄 - IMF报告：吉尔吉斯共和国的真实货币政策远比政策利率宽松
  \item [R005] 智库/国际机构 | 智库/国际机构 / 全球储备资产：SDR新分配暂缓，但危机触发条件仍存 - IMF：SDR新分配的门已经关上，但钥匙还在桌上
  \item [R006] 智库/国际机构 | 智库/国际机构 / 代币化：混合治理模型成为最可能终局，兼容者将胜出 - IMF报告：Tokenization正在重塑金融基础设施，但真正的赢家是“兼容者”
  \item [R007] 战略咨询 | 战略咨询 / 资管行业：增长幻觉终结，净流入能力成为唯一护城河 - 波士顿咨询：资管行业增长幻觉终结，净流入能力成为唯一护城河
  \item [R008] 战略咨询 | 战略咨询 / AI部署：银行'谨慎'正在变成竞争劣势，组织架构比算法更关键 - 波士顿咨询：营销的终极形态，不是“旅程”而是“原子动作”
  \item [R009] 战略咨询 | 战略咨询 / 物理AI：回报周期从5-7年缩至1-3年，产业决策者进入'必须行动'阶段 - 波士顿咨询：BCG：Physical AI的回报周期已从5-7年缩至1-3年，CEO们还在等什么？
  \item [R010] 战略咨询 | 战略咨询 / AI部署：银行'谨慎'正在变成竞争劣势，组织架构比算法更关键 - 波士顿咨询：营销战役的终结，不是营销的终结
  \item [R011] 战略咨询 | 战略咨询 / AI部署：银行'谨慎'正在变成竞争劣势，组织架构比算法更关键 - 波士顿咨询：下一代个性化决策，关键不在算法而在组织架构
  \item [R012] 战略咨询 | 战略咨询 / AI部署：银行'谨慎'正在变成竞争劣势，组织架构比算法更关键 - 麦肯锡：银行对AI的“谨慎”正在变成一种竞争劣势
\end{itemize}
\section{Disclaimer}
{\scriptsize\color{refgray}
This community is a paid learning and information-sharing community created for general educational purposes only. The membership fee is charged solely for access to community discussions, educational content, organization, and operational costs. It is not an advisory fee, management fee, performance fee, brokerage fee, or compensation for personalized financial, investment, legal, tax, or professional advice.\par
I participate and share information solely as an individual and for educational discussion among community members. I am not acting as a financial adviser, investment adviser, broker, dealer, portfolio manager, legal adviser, tax adviser, accountant, fiduciary, or any other licensed professional adviser. Nothing shared in this community constitutes, and should not be interpreted as, financial advice, investment advice, legal advice, tax advice, a recommendation, solicitation, offer, endorsement, instruction, or guarantee to buy, sell, hold, short, trade, or invest in any security, cryptocurrency, fund, derivative, strategy, or other financial product.\par
All content shared in this community, including messages, discussions, links, files, charts, examples, market commentary, watchlists, AI-generated or AI-polished summaries, and third-party materials, is general, impersonal, and educational in nature. It does not take into account any individual's financial situation, investment objectives, risk tolerance, investment horizon, liquidity needs, tax status, legal restrictions, or suitability requirements. No content should be relied upon as suitable for any specific person, account, portfolio, or financial decision.\par
Information shared in this community may be incomplete, inaccurate, outdated, speculative, AI-generated, AI-edited, or based on third-party internet sources that have not been independently verified. I make no representation or warranty as to the accuracy, completeness, timeliness, reliability, or suitability of any information shared.\par
Financial markets involve substantial risk, including the possible loss of principal. Past performance, historical data, hypothetical examples, simulated results, backtests, personal opinions, or educational examples do not guarantee future results. Each member is solely responsible for conducting their own research, exercising independent judgment, and making their own decisions. Before making any financial, investment, legal, or tax decision, members should consult qualified and licensed professionals.\par
Participation in this community, payment of any membership fee, or communication with me or other members does not create any adviser-client, fiduciary, professional, contractual, agency, partnership, employment, or other advisory relationship. By joining or participating in this community, you acknowledge and agree that you are solely responsible for your own decisions, actions, transactions, profits, losses, risks, and consequences, and that I shall not be liable for any direct, indirect, incidental, consequential, financial, legal, tax, or other loss or damage arising from or related to any content, discussion, information, or material shared in this community.\par
}
\end{document}
